\documentclass[12pt]{article}

\usepackage[margin=1in]{geometry}
\usepackage{setspace}
\usepackage[T1]{fontenc}
\usepackage{amsmath,amssymb,amsthm,mathtools,bm}
\usepackage{booktabs}
\usepackage{array}
\usepackage{enumitem}
\usepackage{hyperref}
\usepackage{natbib}
\usepackage{xcolor}
\usepackage{verbatim}
\usepackage{caption}

\onehalfspacing
\hypersetup{
  colorlinks=true,
  linkcolor=blue,
  citecolor=blue,
  urlcolor=blue
}

\newcommand{\R}{\mathbb{R}}
\newcommand{\E}{\mathbb{E}}
\newcommand{\Var}{\operatorname{Var}}
\newcommand{\N}{\mathcal{N}}
\newcommand{\eps}{\varepsilon}
\newcommand{\thetafull}{\theta}
\newcommand{\given}{\,\middle|\,}
\newcommand{\Sigmay}{\Sigma_y}
\newcommand{\norm}[1]{\left\lVert #1 \right\rVert}
\newcommand{\rom}{\mathrm{ROM1}}
\newcommand{\sep}{\mathrm{SEP}}
\newcommand{\sur}{\mathrm{sur}}
\newcommand{\obs}{\mathrm{obs}}
\newcommand{\llik}{\ell}

\newtheorem{protocol}{Protocol}
\newtheorem{diagnostic}{Diagnostic}

% Lightweight algorithm environment, avoiding extra package dependencies.
\newcounter{algorithm}
\newenvironment{algorithm}[1][]{%
  \refstepcounter{algorithm}%
  \par\medskip\noindent
  \textbf{Algorithm \thealgorithm}\if\relax\detokenize{#1}\relax\else: #1\fi
  \par\smallskip
  \begin{enumerate}[leftmargin=2.2em,itemsep=0.25em]
}{%
  \end{enumerate}
  \par\medskip
}

\title{Online Technical Appendix\\
\large Computational Pipeline for ROM1-Residual Surrogate Estimation}
\author{Matyas Farkas}
\date{May 2026}

\begin{document}
\maketitle

\begin{abstract}
\noindent
This appendix documents the computational pipeline used to estimate nonlinear
DSGE models with a ROM1-residual neural-network surrogate. The goal is
replication and auditability. The appendix defines the likelihood objects,
offline training design, stochastic extended path (SEP) data-generation step,
ROM1-residual surrogate, inversion-filter objective, regime-switching gate,
NUTS-HMC implementation, direct nonlinear validation checks, and artifact
schema. It distinguishes three objects that are intentionally separate in the
paper: a direct nonlinear transition benchmark, a surrogate approximation to
that benchmark, and the profiled inversion objective used for empirical
posterior comparisons.
\end{abstract}

\tableofcontents
\newpage

\section{Scope}
\label{apptech:scope}

This appendix is a technical companion to the paper. It is not a second
empirical paper. Its role is to make the estimation pipeline inspectable enough
that a reader can see which numerical objects are being compared and which
claims each validation exercise supports.

The core computational claim is narrow:
\[
  \text{ROM1} + \widehat{\Delta}
  \quad\text{can approximate}\quad
  \text{SEP}
\]
on explicitly documented state--shock--parameter support, where
\(\widehat{\Delta}\) is a learned residual correction. The empirical posterior
comparisons in the paper use this approximation inside a profiled inversion
objective. They are therefore not marginal-likelihood comparisons unless an
additional change-of-variables or Laplace determinant is supplied. The
validation exercises below are designed around that distinction.

\paragraph{Maintained components.}
The maintained replication package contains three layers.
\begin{enumerate}[leftmargin=2em,itemsep=0.2em]
\item A direct nonlinear solver layer based on stochastic extended path
      \citep{fair1983solution,adjemian2025stochastic}.
\item A ROM1-residual surrogate layer trained on
      \(g^{\sep}_\theta-g^{\rom}_\theta\), not on the full transition.
\item An online inference layer using shock inversion, optional regime
      switching, and NUTS-HMC \citep{hoffman2014nuts}.
\end{enumerate}

\paragraph{Main limitation.}
The full 18-parameter Smets--Wouters--HLT application is not validated by a
full direct-SEP HMC run. The current medium-scale bridge is a finite-support,
one-period known-feature comparison over the investment block. It is strong
evidence that the residual target is learned on the relevant nonlinear channel;
it is not a complete certificate for the full multi-period inversion-filter
posterior.

\section{Notation and Objects}
\label{apptech:notation}

Let \(s_t\in\R^{n_s}\) denote the model state, \(\eps_t\in\R^{n_\eps}\) the
structural shocks, \(y_t\in\R^{n_y}\) the observables, and
\(\theta\in\Theta\subset\R^p\) the estimated parameter vector. The nonlinear
one-step transition is
\[
  (x_t^{\sep},s_t^{\sep})
  =
  g^{\sep}_\theta(s_{t-1},\eps_t),
\]
where \(x_t^{\sep}\) stacks the quantities retained for the surrogate target.
The first-order perturbation approximation is
\[
  (x_t^{\rom},s_t^{\rom})
  =
  g^{\rom}_\theta(s_{t-1},\eps_t).
\]
The surrogate does not approximate \(g^{\sep}_\theta\) directly. It approximates
the residual
\begin{equation}
  \Delta_\theta(s_{t-1},\eps_t)
  =
  x_t^{\sep}-x_t^{\rom}.
\label{eq:tech_residual_target}
\end{equation}
The online predictor is therefore
\begin{equation}
  \widehat{x}_t
  =
  x_t^{\rom}
  +
  \widehat{\Delta}_w(s_{t-1},\eps_t,\theta),
\label{eq:tech_surrogate_predictor}
\end{equation}
with neural-network weights \(w\).

The observation equation used for the inversion objective is
\begin{equation}
  y_t^{\obs}=H x_t + u_t,\qquad u_t\sim\N(0,\Sigmay).
\label{eq:tech_obs_eq}
\end{equation}
In the HLT bridge and hard-ELB Gal\'i checks, \(\Sigmay\) is fixed by the
validation design. In the empirical HLT estimation, \(\Sigmay\) is a calibrated
measurement-error scale, with robustness checks tied to surrogate validation
RMSE.

\section{End-to-End Data Flow}
\label{apptech:dataflow}

\begin{algorithm}[Pipeline]
\item Choose a model, parameter block, parameter support, observables, and
      shock/state design.
\item For each design point, compute a direct SEP transition and the
      corresponding ROM1 transition at the same
      \((s_{t-1},\eps_t,\theta)\).
\item Save a training payload with inputs
      \(X=[s_{t-1};\eps_t;\theta]\), direct outputs \(Y\), ROM1 outputs
      \(Y^{\rom}\), solver residual diagnostics, observable indices, parameter
      names, and git provenance.
\item Train a neural network on \(Y-Y^{\rom}\). In the HLT bridge the reported
      target is observation-only; in the broader estimation workflow the bundle
      can also carry state corrections needed for propagation.
\item Validate the surrogate out of sample by comparing
      \(Y^{\rom}+\widehat{\Delta}\) to \(Y\), and by comparing posterior
      surfaces against direct SEP on finite support.
\item In empirical estimation, recover shocks period by period using the ROM1
      inversion filter, evaluate the surrogate correction on gate periods, use
      Kalman likelihood contributions on calm periods, and sample \(\theta\)
      with NUTS-HMC.
\end{algorithm}

The artifact convention is deliberately simple. Raw payloads are serialized as
Julia `.jls' files under `.local_artifacts/'. Human-readable summaries are
stored as `SUMMARY.md' or `SUPPORT_REPORT.md'. Tables promoted to the paper are
stored as `.tex' fragments under either the paper directory or the relevant
artifact directory. The paper should not report a numerical result unless a
matching raw payload and human-readable summary exist.

\section{Direct SEP Data Generation}
\label{apptech:sep_generation}

The direct benchmark solves the nonlinear model with stochastic extended path.
For each design point, the solver receives the current parameter vector, a
finite shock path, and the previous state. It returns a nonlinear transition
and a residual diagnostic. A sample is accepted for training or validation only
if the relevant quantities are finite and the final SEP residual is below the
stage-specific acceptance tolerance.

\subsection{Input Design}

The HLT training and bridge scripts use three types of parameter designs.
\begin{enumerate}[leftmargin=2em,itemsep=0.2em]
\item \emph{Prior/Sobol designs} for broad surrogate coverage.
\item \emph{Deterministic grids} for support mapping and posterior-surface
      validation.
\item \emph{Trimmed supported grids} after direct SEP reveals infeasible
      high-stress corners.
\end{enumerate}

For the medium-scale HLT bridge, the maintained parameter block is
\[
  (\rho_b,\rho_{qs},\sigma_b,\sigma_{qs})
  \equiv
  (\texttt{crhob},\texttt{crhoqs},\texttt{z\_eb},\texttt{z\_eqs}).
\]
The supported bridge trims the risk-premium shock scale to
\[
  z_{eb}\in[1.20,1.85],
\]
because the support map showed systematic direct-SEP failure at the
high-stress \(z_{eb}=2.50\) edge. This is not imputation. The benchmark support
is redefined to the region where direct nonlinear solutions are available.

\subsection{ROM1 and SEP at Common Inputs}

For every accepted sample \(j\), the dataset stores
\[
  X_j=[s_{j,-};\eps_j;\theta_j],\qquad
  Y_j=g^{\sep}_{\theta_j}(s_{j,-},\eps_j),\qquad
  Y^{\rom}_j=g^{\rom}_{\theta_j}(s_{j,-},\eps_j).
\]
The residual target is then \(Y_j-Y^{\rom}_j\). This common-input design is the
reason the validation RMSE is meaningful: the network is judged against the
direct nonlinear correction at the same state, shock, and parameter vector.

\subsection{Warm Starts and Failure Handling}

Across deterministic grids, the solver uses the previous converged solution as
an initial guess for the next point. The warm start is accepted only when the
previous solve converged. If the previous solve failed or exceeded tolerance,
the next solve falls back to a cold start from the perturbation solution. This
prevents failed trajectories from propagating across the grid.

Nonfinite samples are dropped before training. Failed direct-SEP cells are
recorded in support reports rather than silently replaced. A grid is promoted
to validation only if its support is explicit: either every cell solves or the
accepted support is trimmed and re-run.

\section{ROM1-Residual Surrogate}
\label{apptech:surrogate}

\subsection{Residual Target}

The trained object is a map
\[
  \widehat{\Delta}_w:\R^{n_s+n_\eps+p}\rightarrow\R^{d_o},
\]
where \(d_o\) is either the number of observables or the number of retained
state/observable outputs. The HLT bridge uses observation-only output,
\[
  \widehat{\Delta}_{w,o}
  \approx
  H\{g^{\sep}_\theta(s,\eps)-g^{\rom}_\theta(s,\eps)\},
\]
because the posterior-grid comparison is a known-feature, one-period
observable comparison. The empirical HLT estimation bundle also contains the
state-propagation machinery needed by the inversion filter.

Learning the residual rather than the full transition has two practical
consequences. First, the target has a smaller scale because ROM1 captures the
local dynamics near steady state. Second, a zero residual is an interpretable
ablation: it gives the same gate and inversion architecture with no neural
nonlinear correction.

\subsection{Training Objective}

Let \(D_{\mathrm{train}}\) be the training set. Inputs and targets are
standardized dimension by dimension:
\[
  \widetilde{X}_{j,k}={X_{j,k}-\bar X_k\over s_{X,k}},\qquad
  \widetilde{\Delta}_{j,\ell}=
  {\Delta_{j,\ell}-\bar\Delta_\ell\over s_{\Delta,\ell}}.
\]
The default objective is mean squared error,
\begin{equation}
  \min_w
  {1\over |D_{\mathrm{train}}|}
  \sum_{j\in D_{\mathrm{train}}}
  \norm{
    \widehat{\Delta}_w(\widetilde X_j)
    -
    \widetilde{\Delta}_j
  }^2.
\label{eq:tech_training_loss}
\end{equation}
Training uses AdamW with gradient clipping, cosine learning-rate decay with
warmup, and SiLU activation. The maintained MLP has two hidden layers in the HLT
bridge runs; the utility layer also supports a residual network architecture
for wider training designs.

\subsection{Validation Metrics}

For each output coordinate \(k\), the validation RMSE is
\[
  \mathrm{RMSE}_k^{\sur}
  =
  \left[
    {1\over |D_{\mathrm{val}}|}
    \sum_{j\in D_{\mathrm{val}}}
    \left(
      Y_{j,k}
      -
      Y^{\rom}_{j,k}
      -
      \widehat{\Delta}_{w,k}(X_j)
    \right)^2
  \right]^{1/2}.
\]
The improvement over ROM1 is
\[
  1-{\mathrm{RMSE}^{\sur}_k\over \mathrm{RMSE}^{\rom}_k}.
\]
The paper reports the aggregate RMSE and surface-RMSE diagnostics only after
the validation point is held out from training.

\section{Regime-Switching Inversion Objective}
\label{apptech:inversion}

The empirical HLT posterior uses a profiled inversion objective. The state
dimension and shock sequence are too large to sample jointly with the
18-dimensional parameter vector in the current implementation. The inversion
filter reduces the online sampler to \(\theta\) by recovering shocks period by
period.

\subsection{Shock Recovery}

Given \(s_{t-1}\), \(\theta\), and observed \(y_t^{\obs}\), the ROM1 inversion
step solves
\begin{equation}
  \widehat{\eps}_t(\theta)
  =
  \arg\min_{\eps}
  \left\|
    {y_t^{\obs}-H g^{\rom}_\theta(s_{t-1},\eps)\over \sigma_y}
  \right\|^2
  +
  \left\|
    {\eps\over \sigma_\eps(\theta)}
  \right\|^2 .
\label{eq:tech_rom_inversion}
\end{equation}
The same recovered shock can then be evaluated under ROM1, direct SEP in
validation exercises, or ROM1 plus the surrogate residual. This common-shock
design is essential: it isolates the transition approximation from shock
recovery.

\subsection{Profiled Observation Objective}

On a surrogate-evaluated period, the prediction is
\[
  \widehat y_t^{\sur}(\theta)
  =
  H\left[
    g^{\rom}_\theta(s_{t-1},\widehat{\eps}_t)
    +
    \widehat{\Delta}_w(s_{t-1},\widehat{\eps}_t,\theta)
  \right].
\]
The period contribution is the Gaussian observation objective
\begin{equation}
  \ell_t^{\sur}(\theta)
  =
  -{1\over 2}
  \left[
    (y_t^{\obs}-\widehat y_t^{\sur})'
      \Sigmay^{-1}
    (y_t^{\obs}-\widehat y_t^{\sur})
    +
    \log |2\pi \Sigmay|
  \right].
\label{eq:tech_sur_objective}
\end{equation}
Because \(\widehat{\eps}_t(\theta)\) is profiled out, this is not the exact
marginal likelihood unless the determinant terms associated with the profiling
step are added. The paper therefore refers to the reported values as profiled
log objectives.

\subsection{Gate}

The regime switch is deterministic conditional on a gate-calibration payload.
Let \(m_t\in\{0,1\}\) denote the hard gate, with \(m_t=1\) meaning that the
surrogate correction is active. The gate is built from shock and forecast-error
statistics, padded before and after detected stress periods, and filtered to
avoid isolated one-period switches. With a hard gate,
\[
  \ell_T(\theta)
  =
  \sum_{t:m_t=1} \ell_t^{\sur}(\theta)
  +
  \sum_{t:m_t=0} \ell_t^{K}(\theta),
\]
where \(\ell_t^K\) is the Kalman contribution under ROM1. With a soft gate, the
implementation uses precomputed probabilities \(p_t\in[10^{-4},1-10^{-4}]\)
and combines period contributions by the maintained mixing routine. The q95
gate diagnostic in the paper recalibrates this object to avoid classifying all
quiet-sample periods as stressed.

\section{Prior, Parameter Transform, and HMC}
\label{apptech:hmc}

Each estimated parameter has finite support \([a_i,b_i]\). NUTS operates on an
unconstrained vector \(z\in\R^p\), with
\[
  \theta_i(z_i)
  =
  a_i + (b_i-a_i)\Lambda(z_i),
  \qquad
  \Lambda(z)={1\over 1+\exp(-z)}.
\]
The log posterior in unconstrained coordinates is
\begin{equation}
  \log \pi(z\given y)
  =
  \ell_T(\theta(z))
  +
  \log p_\Theta(\theta(z))
  +
  \sum_{i=1}^p
  \log\{(b_i-a_i)\Lambda(z_i)[1-\Lambda(z_i)]\}.
\label{eq:tech_unconstrained_posterior}
\end{equation}
Gradients are computed by finite differences in the maintained HLT estimation
script. The sampler is AdvancedHMC.jl's multinomial NUTS implementation. The
reported diagnostics are divergent transitions, numerical-error proposals,
acceptance rate, split \(\widehat R\), and rank-normalized effective sample size
\citep{vehtari2021rank}.

\begin{protocol}[HMC comparability]
When comparing direct and surrogate objectives in validation, the following
quantities are held fixed: data, priors, bounded-to-unbounded transform,
initial values, seeds, warmup length, post-warmup draws, finite-difference step,
target acceptance rate, and tree-depth limit. A validation failure is therefore
interpreted as an objective difference or numerical-stability difference, not as
a sampler-design difference.
\end{protocol}

\section{Validation Package}
\label{apptech:validation_package}

The validation package has two complementary roles. The Gal\'i hard-ELB
exercise validates pipeline layers in a small model where direct nonlinear
evaluation is feasible. The reduced HLT bridge validates the residual target in
the medium-scale model used for the main application.

\subsection{Gal\'i Hard-ELB Package}

The Gal\'i package uses the maintained current-repo Gal\'i OBC model and three
observables: output, inflation, and the policy rate. It imposes an adverse
\(\eps_z\) shock block that makes the lower bound bind for multiple periods,
with output and inflation falling during the binding window. The package
contains five checks.

\begin{enumerate}[leftmargin=2em,itemsep=0.2em]
\item A stress-path check verifying that the hard-ELB path is economically
      meaningful.
\item A known-shock residual-grid check over \(\sigma_z\).
\item A one-parameter inversion-grid check using common ROM1-recovered shocks.
\item A one-parameter matched-HMC smoke using direct and surrogate objectives.
\item A two-parameter matched-HMC diagnostic over \((\sigma_z,\sigma_a)\).
\end{enumerate}

The extended two-parameter run is the strongest Gal\'i validation currently
reported. It uses a 24-period actual-floor path, four chains per objective, and
4,000 post-warmup draws per objective. Direct and surrogate 90 percent
intervals overlap for both parameters, both true values are covered, and
posterior mean differences are below one combined-MCSE unit in absolute value.
Probes involving \(\sigma_\nu\) are retained as identification diagnostics:
direct and surrogate objectives agree, but the parameter is pushed to the local
grid boundary.

\subsection{Reduced HLT Bridge}

The reduced HLT bridge targets the investment and risk-premium block
\[
  (\texttt{crhob},\texttt{crhoqs},\texttt{z\_eb},\texttt{z\_eqs}),
\]
with observables
\[
  \texttt{dy},\quad \texttt{dinve},\quad \texttt{labobs},\quad
  \texttt{pinfobs},\quad \texttt{robs}.
\]
The maintained 10-by-10-by-10-by-10 supported grid solves all 10,000 direct SEP
cells. The support report records SEP residual min/median/max equal to
\[
  2.700\times 10^{-13},\qquad
  1.571\times 10^{-9},\qquad
  9.996\times 10^{-6}.
\]
The posterior-grid comparison selects a held-out truth point from the validation
split. It constructs a Gaussian observation scale, computes direct SEP, ROM1,
and surrogate log-posterior surfaces on the same grid, and summarizes marginal
posterior intervals by discrete posterior weights.

The maintained result is:
\[
  \mathrm{RMSE}(\rom,\sep)=0.143574,\qquad
  \mathrm{RMSE}(\sur,\sep)=0.000749395,
\]
and
\[
  \mathrm{RMSE}_{\log p}(\rom,\sep)=88.0719,\qquad
  \mathrm{RMSE}_{\log p}(\sur,\sep)=0.0534526.
\]
Surrogate 90 percent intervals overlap direct SEP intervals for all four bridge
parameters. This validates the residual-learning target on finite HLT support.
It does not validate full multi-period inversion-filter HMC for the
18-parameter empirical posterior.

\section{Acceptance Criteria}
\label{apptech:criteria}

\begin{table}[htbp]
\centering
\caption{Validation Criteria by Layer}
\label{tab:tech_validation_criteria}
\begin{tabular}{p{0.27\textwidth}p{0.62\textwidth}}
\toprule
Layer & Criterion \\
\midrule
SEP support & Direct nonlinear cells are finite and either all solve on the
reported support or failures are explicitly mapped and the support is trimmed. \\
Training payload & \(X\), \(Y\), \(Y^{\rom}\), metadata, observable indices,
parameter names, solver residuals, and git provenance are saved. \\
Surrogate accuracy & Held-out RMSE of \(Y^{\rom}+\widehat{\Delta}\) is reported
against direct SEP; ROM1 baseline RMSE is reported on the same split. \\
Posterior grid & Direct SEP and surrogate posterior intervals overlap for all
reported bridge parameters; surrogate surface RMSE is materially below ROM1
surface RMSE. \\
Matched HMC & Direct and surrogate intervals overlap, posterior mean
differences are within combined Monte Carlo error for the targeted parameters,
and numerical-error proposals are zero or explicitly negligible. \\
Empirical use & Any paper table has a raw payload, human-readable summary, and
provenance note. \\
\bottomrule
\end{tabular}
\end{table}

\section{Artifact Schema and Replication Entry Points}
\label{apptech:artifacts}

\subsection{Main Scripts}

The maintained scripts are organized by layer.
\begin{itemize}[leftmargin=2em,itemsep=0.2em]
\item `scripts/hlt\_sep\_surrogate\_dataset\_generate.jl':
      generates direct SEP/ROM1 training payloads.
\item `scripts/hlt\_sep\_surrogate\_train.jl':
      trains the ROM1-residual surrogate.
\item `scripts/hlt\_bridge\_support\_report.jl':
      converts dataset metadata into finite-support reports.
\item `scripts/hlt\_bridge\_posterior\_grid\_compare.jl':
      compares direct SEP, ROM1, and surrogate posterior-grid surfaces.
\item `scripts/hlt\_bridge\_robustness\_sweep.jl':
      orchestrates the HLT bridge dataset, support report, training, and
      posterior-grid comparison.
\item `scripts/run\_surrogate\_hmc\_advancedhmc.jl':
      runs empirical HLT surrogate NUTS-HMC with inversion filtering and gate
      support.
\item `scripts/gali\_validation\_package.jl':
      audits the maintained Gal\'i hard-ELB validation artifacts.
\item `scripts/gali\_obc\_actual\_floor\_twoparam\_inversion\_grid\_validation.jl':
      runs the two-parameter actual-floor inversion/HMC diagnostic.
\end{itemize}

\subsection{Canonical Artifact Layout}

For an HLT bridge run with identifier \(r\), the expected layout is
\begin{verbatim}
.local_artifacts/hlt_reduced_bridge_validation/r/
  hlt_sep_surrogate_dataset.jls
  SUPPORT_REPORT.md
  hlt_sep_surrogate_trained_r.jls

.local_artifacts/hlt_reduced_bridge_validation/posterior_grid_compare_r/
  hlt_bridge_posterior_grid_comparison.jls
  SUMMARY.md
  comparison_table.tex
  ROBUSTNESS_SWEEP.md
\end{verbatim}

For the Gal\'i package, the consolidated audit is written under
\begin{verbatim}
.local_artifacts/gali_validation_package/
\end{verbatim}
and individual layer artifacts remain under their producing run directories.

\subsection{Replication Commands}

The dense HLT bridge is expensive but has a single orchestration entry point:
\begin{verbatim}
julia --project=. scripts/hlt_bridge_robustness_sweep.jl \
  --stage=full \
  --run-id=<run_id> \
  --param-set=investment_4p_supported \
  --grid=10 \
  --epochs=400 \
  --hidden=256 \
  --hidden2=128 \
  --obs-sigma-scale=1.0 \
  --dgp-noise-scale=0.25
\end{verbatim}

The Gal\'i validation package audit is:
\begin{verbatim}
julia --project=. scripts/gali_validation_package.jl \
  --run-id=<run_id> \
  --strict=true
\end{verbatim}

The empirical HLT surrogate chain is run through:
\begin{verbatim}
julia --project=. scripts/run_surrogate_hmc_advancedhmc.jl \
  --surrogate=<surrogate_bundle.jls> \
  --data=<data_payload.jls> \
  --gate-calibration=<gate_payload.jls> \
  --gate-mode=soft \
  --out=<chain_payload.jls>
\end{verbatim}
The linear+gate ablation uses the same command with
`--disable-nn-correction=true'.

\section{Interpretation Rules}
\label{apptech:interpretation}

The pipeline deliberately separates approximation accuracy from economic
interpretation.
\begin{enumerate}[leftmargin=2em,itemsep=0.25em]
\item A low residual RMSE means the network tracks direct SEP on the audited
      support. It does not by itself prove posterior correctness away from that
      support.
\item A passing Gal\'i two-parameter HMC diagnostic validates the
      ROM1-residual, inversion, and matched-HMC layers in a hard-ELB small
      model. It does not establish that the 18-parameter HLT posterior has been
      globally explored.
\item The HLT 10,000-cell bridge validates the residual target in the
      medium-scale investment block. It does not supply a full multi-period
      direct-SEP HMC benchmark.
\item The empirical profiled objective is comparable across linear, linear+gate,
      and surrogate runs that use the same inversion architecture. It should not
      be described as a Bayes factor.
\item Warm-started HMC diagnostics certify sampling within the reported basin.
      They do not rule out other posterior basins unless a separate mode search
      establishes their mass.
\end{enumerate}

\section{Reproducibility Checklist}
\label{apptech:checklist}

Before a result is promoted to the paper, the replication package should contain
the following.
\begin{enumerate}[leftmargin=2em,itemsep=0.25em]
\item The exact command or orchestration script.
\item The git commit recorded in the payload.
\item A raw `.jls' payload containing draws, grids, objectives, or predictions.
\item A human-readable `SUMMARY.md' or support report.
\item A LaTeX table fragment if the result appears in a table.
\item A statement of support: full grid support, trimmed support, or mapped
      failure region.
\item A statement of scope: direct likelihood, profiled inversion objective,
      one-period known-feature comparison, or full HMC comparison.
\end{enumerate}

\section{Conclusion}
\label{apptech:conclusion}

The maintained pipeline is designed to be auditable rather than opaque. The
surrogate is trained on an observable residual against direct nonlinear
solutions; validation compares ROM1, direct SEP, and the residual surrogate on
common inputs; inference uses a profiled inversion objective whose limitations
are explicit; and numerical results promoted to the paper require raw artifacts.
This structure is the basis for the paper's methodological claim. The remaining
gap for a stronger submission is not hidden implementation detail but scale:
the full multi-period direct-SEP HMC benchmark for the 18-parameter HLT
posterior is still beyond the currently reported validation package.

\bibliographystyle{plainnat}
\bibliography{SurrogateNN_paper}

\end{document}
